\documentclass[12pt,a4paper]{article}
\usepackage[utf8]{inputenc}
\usepackage[spanish,es-noquoting]{babel}
\usepackage[T1]{fontenc}
\usepackage{lmodern}
\usepackage{geometry}
\geometry{top=2cm,bottom=2cm,left=2.5cm,right=2.5cm}
\usepackage{xcolor}
\usepackage{listings}
\usepackage{fancyhdr}
\usepackage{graphicx}
\usepackage{tikz}
\usetikzlibrary{positioning}
\usepackage[most]{tcolorbox}
\newtcolorbox{ejerciciografico}{colback=white,colframe=black!65,boxrule=0.7pt,arc=2pt,left=8pt,right=8pt,top=7pt,bottom=7pt}
\graphicspath{{../../assets/}}
\lstset{language=C,basicstyle=\ttfamily\small,breaklines=true,frame=single,showstringspaces=false}
\setlength{\headheight}{15pt}
\pagestyle{fancy}
\fancyhf{}
\lhead{Monitorías Sistemas Operativos 2026-II}
\rhead{Capítulo 1: Procesos}
\cfoot{\thepage}
\title{\textbf{Taller de Lógica: Árboles y Creación de Procesos}}
\author{Monitorías Sistemas Operativos}
\date{\today}
\begin{document}
\maketitle
\section*{Bloque 1: De especificaciones a código}
\begin{ejerciciografico}
\subsection*{Ejercicio 1}
\begin{center}
\begin{tikzpicture}[node distance=1cm, every node/.style={draw,circle,inner sep=2.2pt,fill=black}]
\node (r) {};
\foreach \i in {1,...,5} {\node (h\i) [below=of r,xshift=\i*1.3cm-3.9cm] {}; \draw (r)--(h\i);}
\node (n1) [below=of h1] {}; \node (n3) [below=of h3] {};
\draw (h1)--(n1); \draw (h3)--(n3);
\end{tikzpicture}
\end{center}
\end{ejerciciografico}
\begin{ejerciciografico}
\subsection*{Ejercicio 2}
\begin{center}
\begin{tikzpicture}[node distance=.85cm, every node/.style={draw,circle,inner sep=2.2pt,fill=black}]
\node (r) {}; \node (a) [below left=of r] {}; \node (b) [below=of r] {}; \node (c) [below right=of r] {};
\draw (r)--(a) (r)--(b) (r)--(c);
\node (a1) [below left=of a] {}; \node (a2) [below right=of a] {}; \draw (a)--(a1) (a)--(a2);
\node (b1) [below=of b] {}; \node (b2) [below=of b1] {}; \draw (b)--(b1) (b1)--(b2);
\end{tikzpicture}
\end{center}
\end{ejerciciografico}
\begin{ejerciciografico}
\subsection*{Ejercicio 3}
\begin{center}
\begin{tikzpicture}[node distance=.85cm, every node/.style={draw,circle,inner sep=2.2pt,fill=black}]
\node (r) {}; \node (a) [below=of r] {}; \node (b) [below=of a] {}; \node (c) [below=of b] {};
\node (l) [right=of a] {}; \draw (r)--(a)--(b)--(c); \draw (a)--(l);
\end{tikzpicture}
\end{center}
\end{ejerciciografico}
\begin{ejerciciografico}
\subsection*{Ejercicio 4}
\begin{center}
\begin{tikzpicture}[node distance=.85cm, every node/.style={draw,circle,inner sep=2.2pt,fill=black}]
\node (r) {};
\foreach \i in {0,...,3} {\node (h\i) [below=of r,xshift=\i*1.8cm-2.7cm] {}; \draw (r)--(h\i);}
\foreach \i in {0,2} {\node (a\i) [below left=of h\i] {}; \node (b\i) [below right=of h\i] {}; \draw (h\i)--(a\i) (h\i)--(b\i);}
\end{tikzpicture}
\end{center}
\end{ejerciciografico}
\begin{ejerciciografico}
\subsection*{Ejercicio 5}
\begin{center}
\begin{tikzpicture}[node distance=.75cm, every node/.style={draw,circle,inner sep=2.2pt,fill=black}]
\node (r) {}; \node (a) [below=of r] {}; \node (b) [below=of a] {}; \node (c) [below=of b] {}; \node (d) [below=of c] {};
\draw (r)--(a)--(b)--(c)--(d);
\end{tikzpicture}
\end{center}
\end{ejerciciografico}
\section*{Bloque 2: De código a árbol}
\textit{Para cada fragmento, determine el número total de procesos, dibuje el árbol e identifique cuáles procesos ejecutan cada \texttt{printf}, si existe.}
\subsection*{Ejercicio 6}
\begin{lstlisting}
for (int i = 0; i < 2; i++) {
    if (fork() == 0) break;
}
\end{lstlisting}
\subsection*{Ejercicio 7}
\begin{lstlisting}
if (fork() == 0) {
    fork();
} else {
    fork();
    fork();
}
\end{lstlisting}
\subsection*{Ejercicio 8}
\begin{lstlisting}
pid_t p = fork();
if (p != 0) {
    if (fork() == 0) fork();
}
\end{lstlisting}
\subsection*{Ejercicio 9}
\begin{lstlisting}
for (int i = 0; i < 3; i++) {
    if (i == 1 && fork() == 0) break;
}
\end{lstlisting}
\subsection*{Ejercicio 10}
\begin{lstlisting}
void rama(int n) {
    if (n == 0) return;
    if (fork() == 0) rama(n - 1);
}
int main(void) {
    rama(3);
    return 0;
}
\end{lstlisting}
\subsection*{Ejercicio 11}
El siguiente fragmento contiene un error frecuente. Explique su efecto, corríjalo y analice el árbol que produce la versión corregida.
\begin{lstlisting}
pid_t p = fork();
if (p = 0) {
    fork();
}
\end{lstlisting}
\end{document}
